%%%%%%%%%%%%%%%%%%%%%%%%%%%%%%%%%%%%%%%%%%%%%%%%%%%%%%%%%%%%%%%%%%%%%%%%
%% DECLARAÇÃO DAS FIGURAS — fonte única                                %%
%%                                                                     %%
%% Este arquivo é lido DUAS VEZES, e a macro \figura se comporta de um  %%
%% jeito em cada leitura:                                               %%
%%                                                                     %%
%%   relatorio-descritivo.tex  gera a listagem exigida pelos arts.       %%
%%                             26, III e 27, V — "A Figura 1 apresenta  %%
%%                             <descrição>.";                            %%
%%   desenhos.tex              posiciona a imagem com o rótulo          %%
%%                             "Figura N" centralizado abaixo dela.      %%
%%                                                                     %%
%% Por isso a numeração e a ordem das figuras são consistentes por      %%
%% construção: não existem duas listas para manter em sincronia.        %%
%%                                                                     %%
%% Sintaxe:                                                             %%
%%     \figura{<arquivo em figuras/, sem extensão>}{<descrição>}         %%
%%     \figura[<opções de \includegraphics>]{<arquivo>}{<descrição>}     %%
%%                                                                     %%
%% A descrição entra na frase "A Figura N apresenta ...", então escreva-a %%
%% em minúscula e sem ponto final. Art. 27, V — especifique a           %%
%% representação gráfica: vista, corte, esquema de circuito, diagrama em %%
%% bloco, fluxograma, gráfico.                                          %%
%%                                                                     %%
%% O argumento opcional controla o tamanho. Sem ele, a figura ocupa a   %%
%% página inteira; com uma largura menor (ex.: width=0.45\linewidth),   %%
%% várias figuras cabem na mesma página, como permite o art. 37 — desde %%
%% que nitidamente separadas umas das outras.                           %%
%%                                                                     %%
%% Exigências de conteúdo das figuras (art. 39):                        %%
%%   I   isentas de texto, admitidos apenas termos indicativos ("corte  %%
%%       AA", "vapor d'água", "aberto", "fechado") e palavras-chave     %%
%%       indispensáveis à compreensão;                                   %%
%%   II  diagramas em bloco e fluxogramas contêm descrições de          %%
%%       atividades e/ou etapas;                                         %%
%%   III os termos indicativos não podem cobrir linha alguma da figura; %%
%%   IV  os sinais de referência identificam o MESMO elemento em todas  %%
%%       as figuras em que ele apareça;                                  %%
%%   V   ordenadas, preferencialmente, conforme o relatório descritivo. %%
%%                                                                     %%
%% Art. 21 — nada de timbre, logotipo, letreiro, assinatura ou rubrica  %%
%% nas páginas de desenho.                                              %%
%%%%%%%%%%%%%%%%%%%%%%%%%%%%%%%%%%%%%%%%%%%%%%%%%%%%%%%%%%%%%%%%%%%%%%%%

\figura{figura-1-vista-em-corte}{a vista em corte longitudinal do dispositivo
de acionamento (1), com o corpo tubular (2), a membrana elástica (3), a haste
de acionamento (4), a mola de retorno (5), o atuador eletromecânico (6), o
sensor de pressão (7), a unidade de controle (8) e o assento de vedação de
dupla face (9)}

\figura{figura-2-fluxograma}{o fluxograma das etapas do processo de controle de
vazão executado pela unidade de controle (8)}
